% =====================================================================
%  Lecture 08 -- Random Access Channel (RACH)             (Persian edition)
% =====================================================================
\documentclass[11pt,a4paper]{article}
\providecommand\ACNroot{../..}
\input{\ACNroot/templates/lecture-fa}

\ACNset{number}{08}
\ACNset{title}{کانال دسترسی تصادفی (\lr{RACH})}
\ACNset{shorttitle}{دسترسی تصادفی (\lr{RACH})}
\ACNset{titleother}{Random Access Channel (RACH)}
\ACNset{part}{بخش دوم -- دسترسی و هستهٔ نسل چهارم (درس‌های ۵ تا ۱۳)}
\ACNset{time}{۶ تا ۸ ساعت (۳ ساعت مطالعه، ۳ تا ۵ ساعت تمرین)}
\ACNset{level}{متوسط}
\ACNset{prereq}{درس‌های ۶ و ۷؛ \lr{ALOHA} شیاردار؛ نمایی‌های مختلط و همبستگی}
\ACNset{gen}{4}

\begin{document}
\ACNcover

% ---------------------------------------------------------------------
\section{اهداف یادگیری}
پس از مطالعهٔ این درس باید بتوانید:
\begin{objectives}
  \item موقعیت‌هایی را که در آن یک \acro{UE} در \acro{LTE} باید دسترسی تصادفی انجام دهد \textbf{برشمارید} و
        \textbf{توضیح دهید} چرا همگام‌سازی فراسو مسئلهٔ اصلی است.
  \item کانال \acro{PRACH} را \textbf{توصیف کنید}: قالب‌های پیش‌نوا، پیشوند چرخه‌ای، زمان محافظ و بیشینهٔ برد
        سلولی که از آن‌ها نتیجه می‌شود.
  \item دنباله‌های زادوف--چو، ویژگی‌های همبستگی آن‌ها و نحوهٔ تولید ۶۴ پیش‌نوا در هر سلول با جابه‌جایی‌های
        چرخشی را \textbf{توضیح دهید}.
  \item دسترسی تصادفی مبتنی بر رقابت (چهارمرحله‌ای) و بدون رقابت و گونهٔ دومرحله‌ای \acro{NR} را \textbf{دنبال کنید}.
  \item مقادیر پیش‌فرستی زمانی و شعاع سلول قابل پشتیبانی با یک پیکربندی جابه‌جایی چرخشی را \textbf{محاسبه کنید}.
  \item برخورد پیش‌نواها و ازدحام \acro{RACH} را \textbf{مدل کنید} و راهکارهایی مانند مسدودسازی ردهٔ دسترسی و
        عقب‌نشینی را \textbf{ارزیابی کنید}.
\end{objectives}

% ---------------------------------------------------------------------
\section{دانش پیش‌نیاز}
به ساختار قاب و زمان‌بندی فراسوی \acro{LTE} از درس ۷ (به‌ویژه اینکه ارسال‌های \acro{PUSCH} کاربران مختلف باید
هم‌تراز در زمان به \acro{eNB} برسند)، مفهوم پیش‌فرستی زمانی در \acro{GSM} (درس ۱)، مدل \lr{ALOHA} شیاردار برای
دسترسی تصادفی از درس مقدماتی شبکه و تعریف همبستگی متقابل دوره‌ای دو دنبالهٔ مختلط نیاز دارید.

\paragraph{واژه‌های کلیدی} واژگان این درس همراه با معادل انگلیسی آن‌ها:
\begin{keyterms}
دسترسی تصادفی & \lr{Random access (RA)} & نخستین دسترسی فراسوی هماهنگ‌نشدهٔ یک \acro{UE}. \\
کانال فیزیکی دسترسی تصادفی & \lr{PRACH} & منابع زمان--فرکانس رزروشده برای پیش‌نواها. \\
پیش‌نوا & \lr{Preamble} & دنبالهٔ معلومی که \acro{UE} در \lr{Msg1} می‌فرستد. \\
دنبالهٔ زادوف--چو & \lr{Zadoff--Chu (ZC)} & دنباله‌ای با دامنهٔ ثابت و خودهمبستگی چرخشی ایده‌آل. \\
جابه‌جایی چرخشی & \lr{Cyclic shift ($N_{\mathrm{CS}}$)} & جابه‌جایی‌ای که از یک ریشه پیش‌نواهای متعامد می‌سازد. \\
پیش‌فرستی زمانی & \lr{Timing advance (TA)} & پیش‌انداختن فراسو برای جبران تأخیر رفت‌وبرگشت. \\
پاسخ دسترسی تصادفی & \lr{RAR} & \lr{Msg2}: پیش‌فرستی زمانی، مجوز فراسو و هویت موقت. \\
حل رقابت & \lr{Contention resolution} & \lr{Msg4}: مشخص می‌کند کدام \acro{UE} برخوردکننده برنده شده است. \\
عقب‌نشینی & \lr{Backoff} & انتظار تصادفی پیش از تلاش جدید. \\
مسدودسازی ردهٔ دسترسی & \lr{Access class barring (ACB)} & کنترل پذیرش احتمالاتی پیش از دسترسی تصادفی. \\
\end{keyterms}

% ---------------------------------------------------------------------
\section{مقدمه و انگیزه}
دسترسی زمان‌بندی‌شده به دو چیز نیاز دارد که شبکه برای یک \acro{UE} تازه‌وارد ندارد: یک \emph{هویت} برای
نشانی‌دهی آن روی \acro{PDCCH} و \emph{هم‌ترازی زمانی فراسو}. \acro{UE}ای که تازه بیدار شده زمان‌بندی فروسو را
(از سیگنال‌های همگام‌سازی) می‌داند اما فاصلهٔ خود تا \acro{eNB} را نمی‌داند. اگر به‌سادگی روی \acro{PUSCH}
ارسال می‌کرد، نمادهایش به اندازهٔ تأخیر رفت‌وبرگشت دیر می‌رسیدند و تعامد \acro{SC-FDMA} با کاربران دیگر را از
بین می‌بردند.

دسترسی تصادفی رویه‌ای است که هر دو مسئله را هم‌زمان حل می‌کند. \acro{UE} یک سیگنال کوتاه و مقاوم (یک
\emph{پیش‌نوا}) را در منابعی که برای این منظور رزرو شده می‌فرستد؛ \acro{eNB} زمان رسیدن آن را اندازه می‌گیرد،
یک اصلاح زمانی و نخستین مجوز را بازمی‌فرستد و سپس دو طرف این احتمال را حل می‌کنند که چند \acro{UE} پیش‌نوای
یکسانی انتخاب کرده باشند. این رویه به‌طور شگفت‌آوری پرتکرار است -- در هر برقراری تماس از حالت بیکار، در هر تحویل،
پس از شکست پیوند رادیویی و هر زمان که یک \acro{UE} ناهمگام داده داشته باشد -- پس ظرفیت، تأخیر و استحکام آن
اهمیت دارد. با اینترنت اشیای انبوه و دسترسی مبتنی بر پرتو در \lr{5G}، دسترسی تصادفی به یکی از پرتحرک‌ترین
بخش‌های مهندسی واسط رادیویی تبدیل شده است.

% ---------------------------------------------------------------------
\section{مفاهیم بنیادی}

\subsection{چه زمانی دسترسی تصادفی آغاز می‌شود}
در \acro{LTE} لایهٔ \acro{MAC} دسترسی تصادفی را در این موارد انجام می‌دهد \cite{ts36321}:
\begin{itemize}
  \item دسترسی اولیه از \lr{RRC\_IDLE} (برقراری اتصال \acro{RRC})؛
  \item برقراری مجدد اتصال \acro{RRC} پس از شکست پیوند رادیویی؛
  \item تحویل، برای همگام شدن با سلول مقصد؛
  \item رسیدن دادهٔ فروسو برای \acro{UE} متصلی که همگامی فراسو را از دست داده است؛
  \item رسیدن دادهٔ فراسو وقتی \acro{UE} ناهمگام است یا منبع درخواست زمان‌بندی ندارد؛
  \item مکان‌یابی، جایی که پیش‌فرستی زمانی لازم است.
\end{itemize}

\subsection{مبتنی بر رقابت در برابر بدون رقابت}
در دسترسی تصادفی \textbf{مبتنی بر رقابت} (\acro{CBRA}) \acro{UE} یک پیش‌نوا را به‌طور تصادفی از یک مجموعهٔ
مشترک برمی‌گزیند؛ ممکن است دو \acro{UE} یک پیش‌نوا را انتخاب کنند؛ پس یک مرحلهٔ حل رقابت لازم است. در دسترسی
تصادفی \textbf{بدون رقابت} (\acro{CFRA}) \acro{eNB} از پیش یک پیش‌نوای اختصاصی تخصیص می‌دهد (برای تحویل یا رسیدن
دادهٔ فروسو)؛ برخوردی رخ نمی‌دهد و رویه پس از \acro{RAR} پایان می‌یابد.

\subsection{دنباله‌های زادوف--چو}
\begin{definitionbox}{دنبالهٔ زادوف--چو}
برای طول اول فرد $N_{\mathrm{ZC}}$ و شاخص ریشهٔ $u\in\{1,\dots,N_{\mathrm{ZC}}-1\}$،
\begin{equation}\label{eq:zc}
  x_u(n) = \exp\!\left(-\jj\,\frac{\pi\,u\,n\,(n+1)}{N_{\mathrm{ZC}}}\right), \qquad n = 0,1,\dots,N_{\mathrm{ZC}}-1 .
\end{equation}
دامنهٔ آن ثابت است، خودهمبستگی دوره‌ای آن برای هر جابه‌جایی غیرصفر صفر است و همبستگی متقابل دو ریشهٔ متفاوت
اندازهٔ ثابت $1/\sqrt{N_{\mathrm{ZC}}}$ (بهنجارشده) دارد.
\end{definitionbox}
این ویژگی‌ها دنباله‌های \acro{ZC} را به پیش‌نواهای ایده‌آل تبدیل می‌کنند: دامنهٔ ثابت یعنی نسبت توان قله به
میانگین پایین؛ خودهمبستگی صفر یعنی \emph{جابه‌جایی‌های چرخشی} یک ریشه بر هم متعامدند و تأخیر یک پیش‌نوای
دریافتی به‌صورت یک قلهٔ همبستگی تمیز ظاهر می‌شود. \acro{LTE} برای قالب‌های ۰ تا ۳ از $N_{\mathrm{ZC}} = 839$ استفاده
می‌کند \cite{ts36211}؛ اسکریپت \lr{\texttt{scripts/validate\_examples.py}} این ویژگی‌ها را برای $u=25$ و $u=29$ به‌صورت
عددی بررسی می‌کند.

% ---------------------------------------------------------------------
\section{توضیحات فنی تفصیلی}

\subsection{کانال \lr{PRACH}}
\acro{PRACH} شش بلوک منبع (\lr{1.08\,MHz}) را در فراسو و در زیرقاب‌هایی که \emph{شاخص پیکربندی \lr{PRACH}} در
\lr{SIB2} اعلام می‌کند اشغال می‌کند. درون این باند، پیش‌نوا از فاصلهٔ زیرحامل \lr{1.25\,kHz} استفاده می‌کند؛ بنابراین
۸۳۹ زیرحامل آن در حدود \lr{1.05\,MHz} جا می‌گیرند. هر پیش‌نوا شامل یک پیشوند چرخه‌ای، دنبالهٔ \acro{ZC}
(۸۰۰ میکروثانیه یا ۱۶۰۰ میکروثانیه با تکرار) و یک زمان محافظ ضمنی است:
\begin{table}[htbp]
\centering\small
\caption{قالب‌های پیش‌نوای \acro{LTE} در \acro{FDD} \cite{ts36211}. زمان محافظ تأخیر رفت‌وبرگشت و در نتیجه برد سلول
$d_{\max} = c\,T_{\mathrm{GT}}/2$ را محدود می‌کند.}\label{tab:formats}
\rowcolors{2}{acnmist}{white}
\begin{tabular}{@{}cccccc@{}}
\toprule
\rowcolor{white}قالب & زیرقاب‌ها & پیشوند [$\mu$s] & دنباله [$\mu$s] & محافظ [$\mu$s] & برد [\lr{km}] \\
\midrule
0 & 1 & 103٫1 & 800  & 96٫9  & 14٫5 \\
1 & 2 & 684٫4 & 800  & 515٫6 & 77٫3 \\
2 & 2 & 203٫1 & 1600 & 196٫9 & 29٫5 \\
3 & 3 & 684٫4 & 1600 & 715٫6 & 107٫3 \\
\bottomrule
\end{tabular}
\end{table}
قالب‌هایی با پیشوند و محافظ طولانی‌تر به سلول‌های بزرگ خدمت می‌دهند؛ قالب‌هایی با دنبالهٔ تکراری انرژی بیشتری برای
سلول‌های محدود به پوشش جمع می‌کنند. قالب ۴ (فقط \acro{TDD}) از دنبالهٔ کوتاهی به طول ۱۳۹ در زیرقاب ویژه استفاده می‌کند.

\subsection{۶۴ پیش‌نوا در هر سلول}
هر سلول ۶۴ پیش‌نوا ارائه می‌دهد. این پیش‌نواها از یک یا چند ریشهٔ \acro{ZC} با جابه‌جایی‌های چرخشی $N_{\mathrm{CS}}$
نمونه‌ای تولید می‌شوند: یک ریشه $\floor{N_{\mathrm{ZC}}/N_{\mathrm{CS}}}$ پیش‌نوا می‌دهد و ریشه‌های بعدی (به ترتیب منطقی
پخش‌شده) تا رسیدن به ۶۴ به کار می‌روند. جابه‌جایی باید از بیشینهٔ تأخیر رفت‌وبرگشت به‌علاوهٔ گسترش تأخیر بیشتر
باشد؛ در غیر این صورت پیش‌نوای یک \acro{UE} دور به‌عنوان جابه‌جایی همسایه آشکار می‌شود (\cref{sec:ncs}). ۶۴ پیش‌نوا
میان گروه‌های مبتنی بر رقابت (\lr{A} و \lr{B}؛ دومی نشان‌دهندهٔ \lr{Msg3} بزرگ‌تر است) و مجموعه‌ای رزروشده برای
استفادهٔ بدون رقابت تقسیم می‌شوند.

\subsection{دسترسی تصادفی مبتنی بر رقابت (چهار مرحله)}
\begin{enumerate}
  \item \textbf{\lr{Msg1} -- پیش‌نوا.} \acro{UE} یک پیش‌نوا را به‌طور تصادفی از گروه خود انتخاب و در فرصت
        \acro{PRACH} بعدی با تخمین توان حلقهٔ باز از روی افت مسیر فروسو ارسال می‌کند.
  \item \textbf{\lr{Msg2} -- پاسخ دسترسی تصادفی (\lr{RAR}).} \acro{eNB} در یک پنجرهٔ پاسخ، روی \acro{PDSCH} و با
        نشانی \lr{RA-RNTI} برگرفته از فرصت \acro{PRACH}، شاخص پیش‌نوای آشکارشده (\acro{RAPID})، یک فرمان پیش‌فرستی
        زمانی (۱۱ بیت)، یک مجوز فراسو برای \lr{Msg3} و یک \lr{C-RNTI} موقت را می‌فرستد. اگر \acro{RAR} نرسد،
        \acro{UE} توان خود را \emph{افزایش پله‌ای} می‌دهد و دوباره تلاش می‌کند.
  \item \textbf{\lr{Msg3} -- ارسال زمان‌بندی‌شده.} \acro{UE} پیش‌فرستی زمانی را اعمال و هویت و درخواست خود را روی
        \acro{PUSCH} با \acro{HARQ} ارسال می‌کند -- برای دسترسی اولیه \msg{RRC Connection Request} همراه با
        \lr{S-TMSI} یا یک مقدار تصادفی.
  \item \textbf{\lr{Msg4} -- حل رقابت.} \acro{eNB} هویتی را که آشکار کرده بازتاب می‌دهد. \acro{UE}ای که هویتش
        منطبق است برنده شده و \lr{C-RNTI} موقت به \lr{C-RNTI} آن تبدیل می‌شود. بقیه عقب‌نشینی را آغاز و دوباره تلاش می‌کنند.
\end{enumerate}
دو \acro{UE} که در یک فرصت پیش‌نوای یکسانی انتخاب کرده‌اند هر دو \acro{RAR} یکسانی دریافت و هر دو \lr{Msg3} را
روی منابع یکسان ارسال می‌کنند؛ معمولاً تنها یکی از \lr{Msg3}ها آشکار می‌شود (اثر تسخیر)؛ به همین دلیل
\lr{Msg4} باید هویت را بازتاب دهد.

\subsection{دسترسی تصادفی بدون رقابت}
برای تحویل، سلول مقصد یک پیش‌نوا رزرو و آن را در فرمان تحویل قرار می‌دهد؛ برای رسیدن دادهٔ فروسو، \acro{eNB}
آن را روی \acro{PDCCH} اعلام می‌کند. \acro{UE} پیش‌نوای اختصاصی را می‌فرستد، \acro{RAR} را همراه با پیش‌فرستی
زمانی و مجوز دریافت می‌کند و رویه پایان می‌یابد -- درس ۴ برای مرحلهٔ دسترسی در تحویل \acro{X2} به همین تکیه داشت.

\subsection{ازدحام \lr{RACH} و کنترل آن}
وقتی دستگاه‌های زیادی هم‌زمان دسترسی می‌کنند -- پس از قطع برق، در یک ورزشگاه، یا وقتی هزاران کنتور هم‌زمان
گزارش می‌دهند -- برخوردها چند برابر می‌شوند، تلاش‌های مجدد بیشتر چند برابر می‌شوند و \acro{PRACH} فرو می‌ریزد.
\acro{LTE} چند راهکار ارائه می‌دهد:
\begin{itemize}
  \item \textbf{مسدودسازی ردهٔ دسترسی (\acro{ACB})}: \lr{SIB2} یک ضریب مسدودسازی $a$ و یک زمان پخش می‌کند؛ پیش
        از دسترسی، \acro{UE} یک عدد یکنواخت می‌کشد و تنها اگر از $a$ کمتر باشد ادامه می‌دهد، وگرنه صبر می‌کند.
        مسدودسازی دسترسی گسترده (\acro{EAB}) دستگاه‌های تحمل‌پذیر در برابر تأخیر را هدف می‌گیرد.
  \item \textbf{نشانگر عقب‌نشینی} در \acro{RAR}: بازندگان زمانی تصادفی تا بیشینهٔ اعلام‌شده صبر می‌کنند.
  \item \textbf{منابع \acro{PRACH} بیشتر}: فرصت‌های \acro{PRACH} اضافی در هر قاب، به بهای ظرفیت فراسو.
  \item \textbf{منابع جداگانه برای اینترنت اشیا}: \acro{NB-IoT} و \acro{LTE-M} کانال‌های دسترسی تصادفی خود را تعریف می‌کنند.
\end{itemize}

\subsection{دسترسی تصادفی در \lr{5G NR}}
\acro{NR} ساختار چهارمرحله‌ای را حفظ می‌کند اما آن را با پرتوها و تأخیر کم سازگار می‌کند \cite{ts38300}: هر بلوک
سیگنال همگام‌سازی (\acro{SSB}، یکی برای هر پرتو) با فرصت‌های \acro{PRACH} مشخصی مرتبط است؛ بنابراین پیش‌نوا به
\acro{gNB} می‌گوید \acro{UE} کدام پرتوی فروسو را بهتر می‌بیند؛ پیش‌نواهایی به طول ۸۳۹ یا ۱۳۹ با چند فاصلهٔ زیرحامل
از سلول‌های روستایی دوربرد تا موج میلی‌متری را پشتیبانی می‌کنند؛ و نسخهٔ ۱۶ \textbf{\acro{RACH} دومرحله‌ای} را
می‌افزاید که در آن \lr{MsgA} پیش‌نوا و محمولهٔ \lr{Msg3} را و \lr{MsgB} پاسخ دسترسی تصادفی و حل رقابت را ترکیب
می‌کند و رفت‌وبرگشت‌ها را نصف می‌کند.

% ---------------------------------------------------------------------
\section{مدل‌های ریاضی}

\subsection{پیش‌فرستی زمانی}\label{sec:ta}
\paragraph{مسئله} \acro{LTE} یک \acro{UE} را با چه دقتی هم‌تراز می‌کند و \acro{UE} تا چه فاصله‌ای می‌تواند باشد؟
\paragraph{مدل} \acro{RAR} مقدار $T_A\in\{0,\dots,1282\}$ را حمل می‌کند؛ \acro{UE} فراسوی خود را به اندازهٔ
$N_{\mathrm{TA}} = 16\,T_A\,T_\mathrm{s}$ با $T_\mathrm{s} = 1/(15000\times2048)$\,\lr{s} پیش می‌اندازد \cite{ts36213}.
\paragraph{استخراج} هر گام متناظر با زمان رفت‌وبرگشت $16T_\mathrm{s} = \qty{0.521}{\micro\second}$ و بنابراین فاصلهٔ یک‌طرفهٔ زیر است:
\begin{equation}\label{eq:tastep}
  \Delta d = \frac{c \cdot 16\,T_\mathrm{s}}{2} = \frac{2.998\times10^8\times0.521\times10^{-6}}{2} \approx \qty{78.1}{\metre},
  \qquad d_{\max} = 1282\,\Delta d \approx \qty{100}{\kilo\metre}.
\end{equation}
\paragraph{تفسیر} گام کسری از پیشوند چرخه‌ای عادی (\lr{4.7\,$\mu$s}) است؛ پس ناهم‌ترازی باقی‌مانده را پیشوند
جذب می‌کند. مقایسه کنید با گام ۵۵۳ متری \acro{GSM} (درس ۱): گام ریزتر بازتاب نماد بسیار کوتاه‌تر \acro{LTE} و
نیاز به نگه داشتن همهٔ کاربران درون پیشوند است. پس از دسترسی، فرمان‌های پیش‌فرستی \acro{MAC} (۶ بیتی، نسبی)
هم‌ترازی را با حرکت \acro{UE} حفظ می‌کنند.

\subsection{جابه‌جایی چرخشی و شعاع سلول}\label{sec:ncs}
\paragraph{مدل} مدت دنباله $T_{\mathrm{SEQ}} = 800\,\mu$s شامل $N_{\mathrm{ZC}} = 839$ نمونه است؛ پس هر نمونه
$T_{\mathrm{SEQ}}/N_{\mathrm{ZC}} = 0.954\,\mu$s طول دارد. \acro{UE}ای در فاصلهٔ $r$ پیش‌نوای خود را به اندازهٔ تأخیر
رفت‌وبرگشت $2r/c$ جابه‌جا می‌کند؛ با گسترش تأخیر چندمسیره $\tau_{\mathrm{ds}}$، انرژی دریافتی بازهٔ
$2r/c + \tau_{\mathrm{ds}}$ را اشغال می‌کند.
\paragraph{شرط} برای ابهام نداشتن، این مقدار باید درون یک ناحیهٔ جابه‌جایی چرخشی جا شود:
\begin{equation}\label{eq:ncs}
  N_{\mathrm{CS}}\,\frac{T_{\mathrm{SEQ}}}{N_{\mathrm{ZC}}} \;\ge\; \frac{2r}{c} + \tau_{\mathrm{ds}}
  \quad\Longrightarrow\quad
  r_{\max} = \frac{c}{2}\left(N_{\mathrm{CS}}\,\frac{T_{\mathrm{SEQ}}}{N_{\mathrm{ZC}}} - \tau_{\mathrm{ds}}\right).
\end{equation}
\paragraph{مثال و تفسیر} با $\tau_{\mathrm{ds}} = 5.2\,\mu$s: $N_{\mathrm{CS}} = 13$ مقدار
$r_{\max} \approx \qty{1.1}{\kilo\metre}$ و $\floor{839/13} = 64$ پیش‌نوا از \emph{یک} ریشه می‌دهد؛
$N_{\mathrm{CS}} = 119$ مقدار $r_{\max}\approx\qty{16.2}{\kilo\metre}$ اما تنها ۷ پیش‌نوا برای هر ریشه می‌دهد؛ پس
$\ceil{64/7} = 10$ ریشه لازم است. ریشه‌های متفاوت متعامد نیستند (همبستگی متقابل $1/\sqrt{839}\approx0.035$)؛
بنابراین سلول‌های بزرگ هشدار کاذب و تداخل پیش‌نوای میان‌سلولی بیشتری دارند -- مصالحهٔ پوشش/استحکام که برای هر
سلول برنامه‌ریزی می‌شود. این فرمول حاشیهٔ آشکارسازی کوچکی را که در عمل به کار می‌رود نادیده می‌گیرد؛ حد زمان
محافظ \cref{tab:formats} نیز باید برقرار باشد.

\subsection{احتمال برخورد}\label{sec:collision}
\paragraph{مدل} $n$ \acro{UE} در یک فرصت \acro{PRACH} رقابت می‌کنند و به‌طور یکنواخت از میان $M$ پیش‌نوای رقابتی
انتخاب می‌کنند (مثلاً $M = 54$ وقتی ۱۰ پیش‌نوا از ۶۴ برای \acro{CFRA} رزرو شده‌اند). یک \acro{UE} موفق است اگر هیچ
\acro{UE} دیگری پیش‌نوای آن را انتخاب نکرده باشد (با چشم‌پوشی از خطاهای آشکارسازی و تسخیر \lr{Msg3}).
\paragraph{استخراج} احتمال اینکه یک \acro{UE} معین روی پیش‌نوای خود تنها باشد برابر است با
\begin{equation}\label{eq:psucc}
  P_\mathrm{s}(n) = \left(1-\frac{1}{M}\right)^{n-1},
\end{equation}
و تعداد مورد انتظار \acro{UE}های موفق در آن فرصت
\begin{equation}\label{eq:nsucc}
  S(n) = n\left(1-\frac{1}{M}\right)^{n-1}.
\end{equation}
با پیوسته در نظر گرفتن $n$، $\mathrm{d}S/\mathrm{d}n = 0$ نتیجه می‌دهد $n^\ast = -1/\ln(1-1/M)\approx M$، که در آن
$S(n^\ast)\approx M/\e$ -- دقیقاً بازده \lr{ALOHA} شیاردار، یعنی $1/\e\approx37\%$ به ازای هر پیش‌نوا.
\paragraph{تفسیر} برای $M = 54$: ۱۰ رقیب هر یک با احتمال ۸۴٫۵٪ موفق‌اند؛ ۵۴ رقیب حدود ۲۰ موفقیت (۳۷٪) می‌دهند؛
۱۵۰ رقیب تنها ۹٫۳ موفقیت -- \emph{کمتر} از ۵۴ رقیب (\cref{fig:collision}). فراتر از نقطهٔ بهینه، بار بیشتر یعنی
گذردهی کمتر، و \acro{UE}های ناموفق دوباره تلاش می‌کنند و بار را بیشتر بالا می‌برند: بازخورد مثبتی که به فروپاشی
می‌انجامد مگر اینکه دسترسی محدود شود (\lr{ACB}، عقب‌نشینی). مدل افزایش پله‌ای توان، تسخیر و فرصت‌های متعدد
\acro{PRACH} در هر قاب را نادیده می‌گیرد.

\begin{acnfigure}
  \acnfig{lecture-08/collision}
  \caption{تعداد مورد انتظار \acro{UE}های موفق در هر فرصت \acro{PRACH} (محور چپ، \cref{eq:nsucc}) و احتمال موفقیت
  هر \acro{UE} (محور راست، \cref{eq:psucc}) برای $M = 54$ پیش‌نوای رقابتی.}
  \label{fig:collision}
\end{acnfigure}

\subsection{تأخیر دسترسی}
اگر هر تلاش به‌طور مستقل با احتمال $P_\mathrm{s}$ موفق شود، تعداد تلاش‌ها هندسی با میانگین $1/P_\mathrm{s}$ است و
میانگین تأخیر دسترسی برابر است با
\begin{equation}\label{eq:delay}
  \bar D_{\mathrm{RA}} \approx T_{\mathrm{RA}} + \Bigl(\frac{1}{P_\mathrm{s}}-1\Bigr)\bigl(T_{\mathrm{fail}} + \bar T_{\mathrm{BO}}\bigr),
\end{equation}
که $T_{\mathrm{RA}}$ مدت یک مبادلهٔ چهارمرحله‌ای موفق، $T_{\mathrm{fail}}$ زمان تشخیص شکست (پایان پنجرهٔ \acro{RAR} یا
زمان‌سنج حل رقابت) و $\bar T_{\mathrm{BO}}$ میانگین عقب‌نشینی است (\cref{exr:delay}).

% ---------------------------------------------------------------------
\section{معماری شبکه و نمودارها}
دسترسی تصادفی رویه‌ای صرفاً رادیویی میان \acro{UE} و \acro{eNB} (\acro{MAC} و \acro{PHY}) است، اما نتیجهٔ آن به
لایه‌های بالاتر خوراک می‌دهد: \lr{Msg3} نخستین پیام \acro{RRC} را حمل می‌کند و \lr{Msg4} (در دسترسی اولیه)
پیام \msg{RRC Connection Setup} را، پس از آن پیام \acro{NAS} \msg{Service Request} یا \msg{Attach Request} به
\acro{MME} می‌رسد (درس ۶). \Cref{fig:prach} پیش‌نوای قالب ۰ و \cref{fig:cbra,fig:cfra,fig:twostep} رویه‌ها را نشان می‌دهند.

\begin{acnfigure}
  \resizebox{\linewidth}{!}{\acnfig{lecture-08/prach-format}}
  \caption{پیش‌نوای قالب ۰. پیشوند چرخه‌ای تأخیر رفت‌وبرگشت (\acro{RTD}) و گسترش تأخیر را جذب می‌کند؛ زمان محافظ
  نمی‌گذارد پیش‌نوای دیررس با زیرقاب بعدی هم‌پوشانی کند.}
  \label{fig:prach}
\end{acnfigure}

% ---------------------------------------------------------------------
\section{رویه‌های گام‌به‌گام}
\begin{procedurebox}{دسترسی تصادفی مبتنی بر رقابت برای دسترسی اولیه}\label{proc:cbra}
\begin{steps}
  \item خواندن \lr{SIB2}: شاخص پیکربندی \acro{PRACH}، دنبالهٔ ریشه، $N_{\mathrm{CS}}$، گروه‌های پیش‌نوا، پارامترهای
        توان، پنجرهٔ \acro{RAR}، زمان‌سنج حل رقابت و پارامترهای \acro{ACB}.
  \item گذر از بررسی مسدودسازی دسترسی (در صورت پیکربندی)، انتخاب گروه \lr{A} یا \lr{B}، انتخاب تصادفی یک پیش‌نوا
        و ارسال آن با توان اولیهٔ $P = \min\{P_{\max},\,P_{\mathrm{target}} + \PL\}$.
  \item پایش \acro{PDCCH} برای \lr{RA-RNTI} در پنجرهٔ \acro{RAR}. اگر \acro{RAR} منطبقی نیامد: افزایش توان به
        اندازهٔ گام پله، اعمال عقب‌نشینی و بازگشت به گام ۲ (تا بیشینهٔ تعداد تلاش‌ها).
  \item با دریافت \acro{RAR}: اعمال پیش‌فرستی زمانی، ذخیرهٔ \lr{C-RNTI} موقت، ارسال \lr{Msg3} روی \acro{PUSCH}
        اعطاشده و آغاز زمان‌سنج حل رقابت.
  \item با دریافت \lr{Msg4} با هویت منطبق: موفقیت، تأیید \lr{C-RNTI} و ادامهٔ \acro{RRC}. پایان زمان‌سنج یا عدم
        انطباق: عقب‌نشینی و تلاش مجدد.
\end{steps}
\end{procedurebox}

\begin{acnfigure}
  \acnfig{lecture-08/cbra}
  \caption{دسترسی تصادفی مبتنی بر رقابت (چهارمرحله‌ای) (\cref{proc:cbra}).}
  \label{fig:cbra}
\end{acnfigure}

\begin{acnfigure}
  \acnfig{lecture-08/cfra}
  \caption{دسترسی تصادفی بدون رقابت، مثلاً در تحویل.}
  \label{fig:cfra}
\end{acnfigure}

\begin{acnfigure}
  \acnfig{lecture-08/two-step}
  \caption{دسترسی تصادفی دومرحله‌ای \acro{NR} (نسخهٔ ۱۶).}
  \label{fig:twostep}
\end{acnfigure}

% ---------------------------------------------------------------------
\section{مثال‌های حل‌شده}

\begin{examplebox}{پیش‌فرستی زمانی از روی تأخیر اندازه‌گیری‌شده}\label{ex:ta}
\textbf{مسئله.} \acro{eNB} تأخیر رفت‌وبرگشت یک پیش‌نوا را \lr{21.4\,$\mu$s} اندازه می‌گیرد. چه $T_A$ای ارسال
می‌شود و \acro{UE} چقدر دور است؟

\textbf{حل.} $T_A = \mathrm{round}(21.4/0.521) = 41$؛ $d \approx 41\times78.1 = \qty{3.2}{\kilo\metre}$
(به‌طور مستقیم: $c\tau/2 = 3.21$\,\lr{km}).
\end{examplebox}

\begin{examplebox}{برنامه‌ریزی $N_{\mathrm{CS}}$}\label{ex:ncs}
\textbf{مسئله.} یک سلول روستایی باید به \acro{UE}ها تا فاصلهٔ ۱۲ کیلومتر با گسترش تأخیر \lr{5.2\,$\mu$s} خدمت دهد.
کوچک‌ترین $N_{\mathrm{CS}}$، تعداد پیش‌نوا در هر ریشه و تعداد ریشه‌ها را بیابید.

\textbf{حل.} طبق \cref{eq:ncs}: $N_{\mathrm{CS}} \ge (2\times12\,000/c + 5.2\times10^{-6})/0.954\times10^{-6} =
(80.1+5.2)/0.954 = 89.4$. مجموعهٔ استاندارد مقادیر $N_{\mathrm{CS}}$ شامل ۹۳ است که $\floor{839/93} = 9$ پیش‌نوا در هر
ریشه و $\ceil{64/9} = 8$ ریشه می‌دهد.
\end{examplebox}

\begin{examplebox}{انفجار رویدادمحور اینترنت اشیا}\label{ex:burst}
\textbf{مسئله.} پس از قطع برق، ۵۰۰۰ کنتور در یک سلول تلاش می‌کنند ظرف ۱ ثانیه دسترسی پیدا کنند. سلول در هر قاب ۱۰
میلی‌ثانیه‌ای یک فرصت \acro{PRACH} (۱۰۰ فرصت در ثانیه) و $M = 54$ پیش‌نوای رقابتی دارد. اگر هر کنتور یک بار در یک
فرصت تصادفی تلاش کند، تعداد موفقیت‌ها را تخمین بزنید و یک ضریب مسدودسازی دسترسی پیشنهاد کنید.

\textbf{حل.} به‌طور میانگین ۵۰ رقیب در هر فرصت: $S = 50\,(53/54)^{49} \approx 50\times0.40 = 20.0$ موفقیت در هر فرصت،
حدود ۲۰۰۰ در ثانیهٔ نخست؛ ۳۰۰۰ بازنده دوباره تلاش می‌کنند و به بار ثانیهٔ بعد افزوده می‌شوند. اگر همهٔ ۵۰۰۰ کنتور
به‌جای آن در ۱۰ فرصت (یکی در هر ۱۰۰ میلی‌ثانیه) می‌افتادند، ۵۰۰ رقیب در هر فرصت وجود داشت و
$S = 500\,(53/54)^{499}\approx 0.04$: تقریباً هیچ‌کس موفق نمی‌شد. ضریب مسدودسازی $a\approx 54/500\approx0.1$ بار مؤثر
را به نقطهٔ بهینهٔ $n\approx M$ بازمی‌گرداند که در آن حدود $M/\e\approx20$ \acro{UE} در هر فرصت موفق می‌شوند.

\textbf{تفسیر.} در یک کانال دسترسی تصادفی پربار، محدودسازی گذردهی را \emph{افزایش} می‌دهد -- درس اصلی \lr{ALOHA}.
\end{examplebox}

% ---------------------------------------------------------------------
\section{مقایسهٔ فناوری‌ها}
\begin{table}[htbp]
\centering\small
\caption{دسترسی تصادفی در نسل‌های مختلف.}\label{tab:compare}
\rowcolors{2}{acnmist}{white}
\begin{tabularx}{\linewidth}{@{}l Y Y Y@{}}
\toprule
\rowcolor{white} & \acro{GSM} & \acro{LTE} & \acro{NR} \\
\midrule
کانال               & بُرست \acro{RACH} روی یک شیار زمانی & \acro{PRACH}، ۶ \acro{RB}، پیش‌نواهای \acro{ZC} & \acro{PRACH} مرتبط با پرتوهای \acro{SSB} \\
حل رقابت            & تخصیص \acro{SDCCH} + هویت           & \lr{Msg3/Msg4} با هویت                 & چهارمرحله‌ای یا دومرحله‌ای (\lr{MsgA/MsgB}) \\
گام / برد \acro{TA} & ۵۵۳ متر / ۳۵ کیلومتر                & ۷۸ متر / ۱۰۰ کیلومتر                   & وابسته به عدد‌شناسی (ریزتر در فاصلهٔ زیرحامل بالاتر) \\
بدون رقابت          & ندارد                               & دارد (تحویل، دادهٔ فروسو)              & دارد، به‌علاوهٔ بازیابی از شکست پرتو \\
کنترل ازدحام        & رده‌های دسترسی                      & \lr{ACB, EAB}، عقب‌نشینی               & کنترل دسترسی یکپارچه \\
\bottomrule
\end{tabularx}
\end{table}

\paragraph{مصالحه‌ها} پیش‌نواهای طولانی‌تر و جابه‌جایی‌های چرخشی بزرگ‌تر برد را افزایش می‌دهند اما تعداد پیش‌نواهای
متعامد در هر ریشه را کاهش می‌دهند؛ فرصت‌های \acro{PRACH} بیشتر برخوردها را کم می‌کنند اما ظرفیت فراسو مصرف می‌کنند؛
\acro{RACH} دومرحله‌ای تأخیر را نصف می‌کند اما به منابع \acro{PUSCH} برای محموله‌هایی که ممکن است برخورد کنند نیاز دارد.

% ---------------------------------------------------------------------
\section{نکته‌های کلیدی}
\begin{takeawaybox}
\begin{itemize}
  \item دسترسی تصادفی هم‌ترازی زمانی فراسو و نخستین هویت/مجوز را فراهم می‌کند؛ برای دسترسی اولیه، برقراری مجدد،
        تحویل و دادهٔ ناهمگام به کار می‌رود.
  \item پیش‌نواها دنباله‌های \acro{ZC} ($N_{\mathrm{ZC}} = 839$) هستند: دامنهٔ ثابت، خودهمبستگی چرخشی صفر؛ ۶۴ عدد در هر
        سلول از جابه‌جایی‌های چرخشی یک یا چند ریشه.
  \item $r_{\max} = \tfrac{c}{2}(N_{\mathrm{CS}}T_{\mathrm{SEQ}}/N_{\mathrm{ZC}} - \tau_{\mathrm{ds}})$؛ زمان محافظ نیز برد را محدود
        می‌کند (قالب ۰: ۱۴٫۵ کیلومتر).
  \item گام \acro{TA} برابر $16T_\mathrm{s} = 0.52\,\mu$s $\Rightarrow$ ۷۸ متر؛ \acro{TA} یازده‌بیتی \acro{RAR} $\Rightarrow$ حدود ۱۰۰ کیلومتر.
  \item \acro{CBRA}: پیش‌نوا $\to$ \acro{RAR} $\to$ \lr{Msg3} $\to$ حل رقابت؛ \acro{CFRA} پس از \acro{RAR} پایان می‌یابد.
  \item موفقیت‌ها در هر فرصت $n(1-1/M)^{n-1}$ در حدود $M/\e$ به قله می‌رسند؛ ازدحام گذردهی را فرو می‌ریزد، پس باید
        محدود کرد (\lr{ACB}، عقب‌نشینی).
  \item \acro{NR} \acro{PRACH} مرتبط با پرتو و \acro{RACH} دومرحله‌ای را می‌افزاید.
\end{itemize}
\end{takeawaybox}

% ---------------------------------------------------------------------
\section{تمرین‌ها}

\begin{exercisebox}{1}{چرا مستقیماً \lr{PUSCH} نه؟}\label{exr:why}
توضیح دهید اگر یک \acro{UE} بیکار نخستین پیام خود را بدون دسترسی تصادفی مستقیماً روی \acro{PUSCH} می‌فرستاد چه
مشکلی پیش می‌آمد. هم زمان‌بندی و هم هویت را در نظر بگیرید.
\end{exercisebox}

\begin{exercisebox}{1}{\lr{CBRA} یا \lr{CFRA}؟}\label{exr:which}
برای هر مورد بگویید از \acro{CBRA} استفاده می‌شود یا \acro{CFRA} و چرا: اتصال اولیه؛ تحویل؛ برقراری مجدد \acro{RRC}
پس از \acro{RLF}؛ رسیدن دادهٔ فروسو برای \acro{UE} متصل ناهمگام.
\end{exercisebox}

\begin{exercisebox}{2}{انتخاب قالب}\label{exr:format}
یک سلول باید ۶۰ کیلومتر را روی زمین هموار پوشش دهد. از نظر زمان محافظ کدام قالب‌های پیش‌نوای \acro{LTE} ممکن‌اند و
کوچک‌ترین $N_{\mathrm{CS}}$ (با $\tau_{\mathrm{ds}} = 5.2\,\mu$s) چیست؟ آیا با دنبالهٔ ۸۰۰ میکروثانیه‌ای ممکن است؟
\end{exercisebox}

\begin{exercisebox}{2}{تأخیر دسترسی}\label{exr:delay}
با $T_{\mathrm{RA}} = 20$\,\lr{ms}، $T_{\mathrm{fail}} = 15$\,\lr{ms} و عقب‌نشینی یکنواخت در $[0,20]$\,\lr{ms}، میانگین
تأخیر دسترسی را از \cref{eq:delay} برای $n = 10$ و $n = 54$ رقیب ($M = 54$) محاسبه کنید.
\end{exercisebox}

\begin{exercisebox}{3}{ویژگی‌های \lr{ZC}}\label{exr:zc}
ثابت کنید وقتی $N_{\mathrm{ZC}}$ اول است، خودهمبستگی دوره‌ای \cref{eq:zc} برای هر جابه‌جایی غیرصفر $\tau$ صفر است.
(راهنمایی: $x_u(n+\tau)x_u^\ast(n)$ را بنویسید و یک سری هندسی در $n$ را جمع بزنید.)
\end{exercisebox}

\begin{exercisebox}{3}{پایداری تلاش‌های مجدد}\label{exr:stab}
\acro{UE}های جدید با میانگین $\lambda$ در هر فرصت \acro{PRACH} می‌رسند؛ هر \acro{UE} ناموفق در فرصت بعدی دوباره تلاش
می‌کند. با استفاده از \cref{eq:nsucc} و تقریب پواسون $S \approx G e^{-G/M}$ برای بار ارائه‌شدهٔ $G$ در هر فرصت، شرط
تعادل را بنویسید، بیشینهٔ $\lambda$ پایدار را بیابید و بحث کنید فراتر از آن چه رخ می‌دهد.
\end{exercisebox}

\subsection*{پاسخ‌ها}
\addcontentsline{toc}{subsection}{پاسخ‌ها}

\begin{solutionbox}{\ref{exr:why}}
زمان‌بندی: \acro{UE} تنها زمان‌بندی فروسو را می‌داند؛ فراسوی آن به اندازهٔ تأخیر رفت‌وبرگشت (تا صدها میکروثانیه) دیر
می‌رسید، با نمادهای \acro{SC-FDMA} کاربران دیگر هم‌پوشانی می‌کرد و تعامد را از بین می‌برد. هویت: \acro{eNB} برای
\acro{UE} هیچ \lr{C-RNTI}ای ندارد و نمی‌تواند مجوزی را به آن نشانی‌دهی کند؛ منابع \acro{PUSCH} زمان‌بندی‌نشده وجود ندارند.
دسترسی تصادفی هر دو را با یک پیش‌نوای مقاوم و تحمل‌پذیر در برابر تأخیر حل می‌کند.
\end{solutionbox}

\begin{solutionbox}{\ref{exr:which}}
اتصال اولیه: \acro{CBRA} (شبکه \acro{UE} را نمی‌شناسد). تحویل: \acro{CFRA} (مقصد پیش‌نوایی رزرو می‌کند و مرحلهٔ رقابت و
زمان صرفه‌جویی می‌شود). برقراری مجدد \acro{RRC}: \acro{CBRA} (برنامه‌ریزی‌نشده). دادهٔ فروسو برای \acro{UE} ناهمگام:
\acro{CFRA} (\acro{eNB} \acro{UE} را می‌شناسد و پیش‌نوای اختصاصی را روی \acro{PDCCH} اعلام می‌کند)، و اگر پیش‌نوای اختصاصی
در دسترس نباشد بازگشت به \acro{CBRA}.
\end{solutionbox}

\begin{solutionbox}{\ref{exr:format}}
زمان محافظ باید رفت‌وبرگشت $2\times60\,000/c = 400\,\mu$s را بپوشاند: قالب‌های ۱ (۵۱۵٫۶ میکروثانیه) و ۳ (۷۱۵٫۶
میکروثانیه) واجد شرایط‌اند؛ ۰ و ۲ نه. \Cref{eq:ncs} مستلزم $N_{\mathrm{CS}} \ge (400+5.2)/0.954 = 425$ نمونه است.
بزرگ‌ترین مقدار غیرصفر مجموعهٔ استاندارد بدون محدودیت ۴۱۹ است ($r_{\max}\approx59$ کیلومتر)، اندکی کمتر از لازم؛ گزینهٔ
باقی‌مانده $N_{\mathrm{CS}} = 0$ (بدون جابه‌جایی چرخشی) است که در آن هر ریشه تنها یک پیش‌نوا می‌دهد و ۶۴ ریشه به کار
می‌رود. این برای هر تأخیری کار می‌کند، اما ۶۴ پیش‌نوا دیگر متعامد نیستند (تنها همبستگی متقابل کم آن‌ها را جدا
می‌کند)؛ پس عملکرد آشکارسازی افت می‌کند -- بهای سلول‌های بسیار بزرگ. طول ۸۰۰ میکروثانیه‌ای دنباله مسئله نیست؛ تکرار
(قالب ۳) تنها به بودجهٔ پیوند کمک می‌کند.
\end{solutionbox}

\begin{solutionbox}{\ref{exr:delay}}
$\bar T_{\mathrm{BO}} = 10$\,\lr{ms}. $n=10$: $P_\mathrm{s} = 0.845$، $1/P_\mathrm{s}-1 = 0.183$،
$\bar D = 20 + 0.183\times25 = 24.6$\,\lr{ms}. $n=54$: $P_\mathrm{s} = 0.371$، $1/P_\mathrm{s}-1 = 1.70$،
$\bar D = 20 + 1.70\times25 = 62.4$\,\lr{ms} (خوش‌بینانه، زیرا تلاش‌های مجدد نیز بار می‌افزایند).
\end{solutionbox}

\begin{solutionbox}{\ref{exr:zc}}
$x_u(n+\tau)x_u^\ast(n) = \exp(-\jj\pi u[(n+\tau)(n+\tau+1) - n(n+1)]/N) =
\exp(-\jj\pi u[\tau^2+\tau]/N)\,\exp(-\jj 2\pi u\tau n/N)$. عامل نخست نسبت به $n$ ثابت است؛ جمع عامل دوم روی
$n = 0,\dots,N-1$ برابر $\sum_n \e^{-\jj2\pi u\tau n/N} = 0$ است مگر اینکه $u\tau \equiv 0 \pmod N$. با $N$ اول و
$1\le u,\tau<N$، حاصل‌ضرب $u\tau$ هرگز مضرب $N$ نیست؛ پس خودهمبستگی برای همهٔ $\tau\neq0$ صفر است. (دوره‌ای بودن
فاز درجه‌دوم به پیمانهٔ $N$ برای $N$ فرد برقرار است.)
\end{solutionbox}

\begin{solutionbox}{\ref{exr:stab}}
در تعادل نرخ موفقیت با نرخ ورود برابر است: $\lambda = G e^{-G/M}$. سمت راست در $G = M$ با مقدار $M/\e$ بیشینه می‌شود؛
پس $\lambda_{\max} = M/\e \approx 19.9$ \acro{UE} در هر فرصت برای $M=54$. برای $\lambda < M/\e$ دو تعادل وجود دارد: تعادل
کم‌بار ($G<M$) پایدار و تعادل پربار ناپایدار است. برای $\lambda > M/\e$ هیچ تعادلی وجود ندارد: صف بی‌کران رشد می‌کند و
موفقیت‌ها به سوی صفر میل می‌کنند -- ناپایداری‌ای که \acro{ACB} و عقب‌نشینی تصادفی برای جلوگیری از آن طراحی شده‌اند.
\end{solutionbox}

% ---------------------------------------------------------------------
\section{مطالعهٔ بیشتر}
رویهٔ دسترسی تصادفی \acro{MAC} در \cite{ts36321}، تولید پیش‌نوا و قالب‌ها در \cite{ts36211}، پیش‌فرستی زمانی در
\cite{ts36213} و \acro{ACB} در مشخصات \acro{RRC} \cite{ts36331} آمده‌اند. دسترسی تصادفی \acro{NR}، از جمله \acro{RACH}
دومرحله‌ای، در \cite{ts38300} خلاصه شده است. \lr{Sesia} و همکاران \cite{sesia2011lte} فصلی مفصل دربارهٔ طراحی
\acro{PRACH} و دنباله‌های \acro{ZC} دارند؛ \lr{Dahlman} و همکاران \cite{dahlman2020nr} رویه‌های \acro{NR} را پوشش می‌دهند.

\begin{latin}
\printbibliography[heading=subbibliography,title={References}]
\end{latin}
\end{document}
